\section{Related work}
\subsection{Current State of the Art and Open Questions}
\label{sec:litrev:sota}

The ALDE benchmark~\cite{yangActiveLearningAssistedDirected} is the most systematic published comparison of
encoding-surrogate-acquisition function combinations to date.
Testing all combinations of four encodings (one-hot, AAIndex/Georgiev, ESM2 site-specific),
four surrogate types (boosting ensemble, GP, DNN ensemble, DKL), and two acquisition functions
(greedy, Thompson sampling) on GB1 and TrpB, it established that:
(i)\ the DNN ensemble with Thompson sampling is the best overall strategy across conditions;
(ii)\ ESM2 site-specific embeddings improve performance when paired with DNN or DKL surrogates;
and (iii)\ non-deep surrogates benefit little or not at all from richer encodings.
These findings are consistent with the earlier observation by Wittmann et al.~\cite{wittmannInformedTrainingSet2021}
that one-hot encoding outperforms more complex representations for non-deep models on GB1.

EVOLVEpro~\cite{jiangRapidProteinEvolution2024} demonstrated a practical active learning system using
ESM2-15B mean-pooled embeddings with a random forest surrogate, achieving measurable fitness
improvements across five diverse protein targets in an iterative wet-lab setting.
EVOLVEpro also reported a notable failure mode: PLM zero-shot fitness scores (log-likelihood
ratios) showed near-zero or negative correlation with experimentally measured activity for
several proteins, highlighting that evolutionary fitness and specific biological function are
distinct quantities.

ESM3~\cite{hayesSimulating500Million} introduced a multimodal generative PLM jointly modelling sequence,
structure, and function tokens.
At 98 billion parameters, the full model is computationally prohibitive for most research
groups; however, smaller ESM C variants (300 million parameters) are more accessible and may
provide embeddings competitive with ESM2 at lower computational cost.
Whether ESM3 embeddings provide material benefits over ESM2 in supervised fitness prediction
relevant to MLDE is not yet established.

Three important questions remain unresolved at the time of writing.
First, essentially all systematic encoding-surrogate benchmarks have been conducted on dense
combinatorial libraries (GB1, TrpB); whether optimal pairings transfer to sparse full-protein
landscapes such as avGFP is unknown.
Second, the delta embedding approach has been proposed as a biologically motivated alternative
to mean pooling but has not been benchmarked.
Third, Precision@K as a primary evaluation metric for active learning in protein engineering
remains underused; it is not clear how strongly encoding-surrogate combinations differ on
this metric versus Spearman $\rho$.
\subsection{Proteins}
Our proteins specifically
\subsection{Thesis contribution}
\label{sec:litrev:rq}

The preceding review motivates the central research question of this thesis:

\begin{quote}
  \emph{XXXX}
\end{quote}

More concretely, we test whether the relative ranking of four encoding strategies
(one-hot mutation features, ESM2 mean-pooled, ESM2 delta, AAIndex physicochemical)
across three surrogate models (random forest, Gaussian process, DNN ensemble) is consistent
between GB1 and avGFP, and whether Spearman $\rho$ and Precision@K give concordant or
divergent rankings.

The work makes four concrete contributions.
A
B
C
D